\documentclass[10pt]{extarticle}
\usepackage[a4paper,margin=1.5cm]{geometry}
\usepackage[T1]{fontenc}
\usepackage{amsmath,amssymb}
\usepackage{graphicx}
\usepackage[export]{adjustbox}
\usepackage{multicol}
\usepackage{enumitem}
\usepackage{booktabs,array,tabularx}
\setlength{\columnsep}{0.7cm}
\setlength{\parindent}{0pt}
\setlength{\parskip}{3pt}
\pagestyle{empty}
\makeatletter
\newcommand{\qfigdef}[4]{\@namedef{qf@file@#1}{#2}\@namedef{qf@w@#1}{#3}\@namedef{qf@wide@#1}{#4}}
\newcommand{\qf@check}[1]{\@ifundefined{qf@file@#1}{\PackageError{score2}{Unknown figure id #1}{}}{}}
\newcommand{\qf@img}[4][]{\includegraphics[width=#2\dimexpr\@nameuse{qf@w@#4}pt\relax,max width=#3,max height=0.45\textheight,#1]{\@nameuse{qf@file@#4}}}
\newif\ifqf@full
\newif\ifqf@boxed
\newsavebox{\qf@box}
\newcommand{\qf@open}{\global\qf@fulltrue\par\global\qf@boxedtrue
  \setbox\qf@box\vbox\bgroup\hsize\textwidth\linewidth\textwidth}
\newcommand{\qf@emit}[1]{\par\egroup\global\qf@boxedfalse
  \ifdim\dimexpr\ht\qf@box+\dp\qf@box\relax<#1\textheight\noindent\usebox{\qf@box}\par\else\unvbox\qf@box\fi}
\newcommand{\qfullsplit}{\ifqf@boxed\qf@emit{0.85}\fi}
\newcommand{\qf@close}{\ifqf@boxed\qf@emit{0.45}\fi\global\qf@fullfalse}
\newcommand{\qfullbegin}{\end{multicols}\qf@open}
\newcommand{\qfullend}{\qf@close\begin{multicols}{2}}
\newcommand{\qfullstart}{\qf@open}
\newcommand{\qfulldone}{\qf@close}
\newcommand{\qfig}[2][1]{\qf@check{#2}%
  \ifqf@full\qf@blk{#1}{#2}\else
  \ifnum\@nameuse{qf@wide@#2}=1 \expandafter\@firstoftwo\else\expandafter\@secondoftwo\fi
  {\end{multicols}\begin{center}\qf@img{#1}{\textwidth}{#2}\end{center}\begin{multicols}{2}}%
  {\qf@blk{#1}{#2}}\fi}
\newcommand{\qf@blk}[2]{\par\smallskip{\centering\qf@img{#1}{\linewidth}{#2}\par}\smallskip}
\newcommand{\qfigi}[2][1]{\qf@check{#2}\qf@img[valign=c]{#1}{\linewidth}{#2}}
\makeatother
\newcommand{\papersec}[1]{\par\bigskip{\centering\Large\bfseries PAPER-#1\par}\medskip}
\newcommand{\partsec}[2]{\par\medskip\noindent{\large\bfseries PART-#1: #2}\par\smallskip}
\newcommand{\sect}[1]{\par\smallskip\noindent\textbf{#1}\par}
\newcommand{\qn}[1]{\par\smallskip\noindent\textbf{#1.}~}
\newcommand{\paraq}[2]{\par\smallskip\noindent\textbf{Paragraph for Questions #1 and #2}\par}
\newenvironment{qoptions}{\begin{itemize}[leftmargin=1.8em,itemsep=1pt,topsep=2pt,label={}]}{\end{itemize}}
\newcommand{\opt}[2]{\item[(#1)] #2}
% figure registry (generated by score2_tools.py finalize)
\qfigdef{p1_phy_s1_q01_1}{Score-II - Test 43 (23 Mar 2025)-img/p1_phy_s1_q01_1.png}{76.0}{0}
\qfigdef{p1_phy_s1_q02_1}{Score-II - Test 43 (23 Mar 2025)-img/p1_phy_s1_q02_1.png}{137.1}{0}
\qfigdef{p1_phy_s1_q04_1}{Score-II - Test 43 (23 Mar 2025)-img/p1_phy_s1_q04_1.png}{106.8}{0}
\qfigdef{p1_phy_s1_q05_1}{Score-II - Test 43 (23 Mar 2025)-img/p1_phy_s1_q05_1.png}{99.3}{0}
\qfigdef{p1_phy_s1_q06_1}{Score-II - Test 43 (23 Mar 2025)-img/p1_phy_s1_q06_1.png}{89.0}{0}
\qfigdef{p1_phy_s1_q07_1}{Score-II - Test 43 (23 Mar 2025)-img/p1_phy_s1_q07_1.png}{111.7}{0}
\qfigdef{p1_phy_s1_q08_1}{Score-II - Test 43 (23 Mar 2025)-img/p1_phy_s1_q08_1.png}{73.2}{0}
\qfigdef{p1_phy_s1_q09_1}{Score-II - Test 43 (23 Mar 2025)-img/p1_phy_s1_q09_1.png}{89.1}{0}
\qfigdef{p1_phy_s1_q10_1}{Score-II - Test 43 (23 Mar 2025)-img/p1_phy_s1_q10_1.png}{95.4}{0}
\qfigdef{p1_phy_s1_q11_1}{Score-II - Test 43 (23 Mar 2025)-img/p1_phy_s1_q11_1.png}{73.2}{0}
\qfigdef{p1_phy_s1_q15_1}{Score-II - Test 43 (23 Mar 2025)-img/p1_phy_s1_q15_1.png}{78.5}{0}
\qfigdef{p1_phy_s2_q01_1}{Score-II - Test 43 (23 Mar 2025)-img/p1_phy_s2_q01_1.png}{42.0}{0}
\qfigdef{p1_phy_s2_q03_1}{Score-II - Test 43 (23 Mar 2025)-img/p1_phy_s2_q03_1.png}{55.1}{0}
\qfigdef{p1_chem_s1_q02_1}{Score-II - Test 43 (23 Mar 2025)-img/p1_chem_s1_q02_1.png}{99.6}{0}
\qfigdef{p1_chem_s1_q03_1}{Score-II - Test 43 (23 Mar 2025)-img/p1_chem_s1_q03_1.png}{96.9}{0}
\qfigdef{p1_chem_s1_q06_1}{Score-II - Test 43 (23 Mar 2025)-img/p1_chem_s1_q06_1.png}{111.8}{0}
\qfigdef{p1_chem_s1_q06_2}{Score-II - Test 43 (23 Mar 2025)-img/p1_chem_s1_q06_2.png}{59.6}{0}
\qfigdef{p1_chem_s1_q06_4}{Score-II - Test 43 (23 Mar 2025)-img/p1_chem_s1_q06_4.png}{65.1}{0}
\qfigdef{p1_chem_s1_q06_3}{Score-II - Test 43 (23 Mar 2025)-img/p1_chem_s1_q06_3.png}{47.5}{0}
\qfigdef{p1_chem_s1_q10_2}{Score-II - Test 43 (23 Mar 2025)-img/p1_chem_s1_q10_2.png}{21.2}{0}
\qfigdef{p1_chem_s1_q10_1}{Score-II - Test 43 (23 Mar 2025)-img/p1_chem_s1_q10_1.png}{36.3}{0}
\qfigdef{p1_chem_s1_q10_3}{Score-II - Test 43 (23 Mar 2025)-img/p1_chem_s1_q10_3.png}{44.9}{0}
\qfigdef{p1_chem_s1_q10_4}{Score-II - Test 43 (23 Mar 2025)-img/p1_chem_s1_q10_4.png}{31.8}{0}
\qfigdef{p1_chem_s1_q13_1}{Score-II - Test 43 (23 Mar 2025)-img/p1_chem_s1_q13_1.png}{118.1}{0}
\qfigdef{p1_chem_s1_q15_1}{Score-II - Test 43 (23 Mar 2025)-img/p1_chem_s1_q15_1.png}{109.9}{0}
\qfigdef{p1_chem_s1_q20_1}{Score-II - Test 43 (23 Mar 2025)-img/p1_chem_s1_q20_1.png}{40.4}{0}
\qfigdef{p1_chem_s1_q20_2}{Score-II - Test 43 (23 Mar 2025)-img/p1_chem_s1_q20_2.png}{31.8}{0}
\qfigdef{p1_chem_s1_q20_3}{Score-II - Test 43 (23 Mar 2025)-img/p1_chem_s1_q20_3.png}{39.6}{0}
\qfigdef{p1_chem_s1_q20_4}{Score-II - Test 43 (23 Mar 2025)-img/p1_chem_s1_q20_4.png}{96.9}{0}
\qfigdef{p1_chem_s1_q20_5}{Score-II - Test 43 (23 Mar 2025)-img/p1_chem_s1_q20_5.png}{34.3}{0}
\qfigdef{p1_chem_s2_q04_1}{Score-II - Test 43 (23 Mar 2025)-img/p1_chem_s2_q04_1.png}{143.8}{0}
\qfigdef{p1_chem_s2_q05_1}{Score-II - Test 43 (23 Mar 2025)-img/p1_chem_s2_q05_1.png}{144.6}{0}
\makeatletter\@beginparpenalty=10000\makeatother
\begin{document}
\begin{center}{\Large\bfseries JEE (Main+Advanced) Score-II --- Test 43 (23 Mar 2025)}\end{center}
\begin{multicols}{2}
\raggedcolumns
\partsec{1}{PHYSICS}
\sect{Section-I}

\qn{1} Find electric field due to uniformly charged hemispherical cap having charge density $\sigma$ at point P :-
\qfig{p1_phy_s1_q01_1}
\begin{qoptions}
\opt{A}{$\dfrac{\sigma}{2\sqrt2\in_0}$}
\opt{B}{$\dfrac{\sigma}{2\sqrt2\in_0}+\dfrac{\sigma}{2\in_0}$}
\opt{C}{$\dfrac{\sigma}{4\in_0}$}
\opt{D}{$\dfrac{\sigma}{2\in_0}$}
\end{qoptions}

\qn{2} A smooth and vertical circular wire frame of radius 2m is fixed inside water as shown. A small bead of specific gravity 0.5 is threaded on the wire and is kept at the origin. If the bead is imparted velocity $V_0$ towards positive $x$ axis, it moves on the wire frame then neglecting effect of viscosity, minimum value of $V_0$ so that it completes vertical circle will be :-
\qfig{p1_phy_s1_q02_1}
\begin{qoptions}
\opt{A}{$5\sqrt2$ m/s}
\opt{B}{$2\sqrt{20}$ m/s}
\opt{C}{$2\sqrt{10}$ m/s}
\opt{D}{$\sqrt{70}$ m/s}
\end{qoptions}

\qn{3} \textbf{Statement-1 :} Experiments revealed that the distance of closest approach to a gold nucleus of an $\alpha$ particle of kinetic energy 5.5 MeV is about $4\times10^{-14}$ m.\\
\textbf{Statement-2 :} Rutherford's calculations take into account both coulombic repulsion force as well as short ranged nuclear force.
\begin{qoptions}
\opt{A}{Statement-1 is true, statement-2 is true and statement-2 is correct explanation for statement-1.}
\opt{B}{Statement-1 is true, statement-2 is true and statement-2 is NOT the correct explanation for statement-1.}
\opt{C}{Statement-1 is true, statement-2 is false.}
\opt{D}{Statement-1 is false, statement-2 is true.}
\end{qoptions}

\qn{4} A system of capacitors and resistors are connected to an ideal battery of emf $V$ as shown in figure. Initially all the capacitors are uncharged and switch is open. Now switch S is closed at $t=0$. (Capacitance of each capacitor is `$C$' and resistance of each resistor is $R$ and all of the symbols have their usual meaning). Current through battery just after switch is closed, is
\qfig{p1_phy_s1_q04_1}
\begin{qoptions}
\opt{A}{$\dfrac7{13}\dfrac VR$}
\opt{B}{$\dfrac{17}{13}\dfrac VR$}
\opt{C}{$\dfrac{2V}R$}
\opt{D}{None of these}
\end{qoptions}

\qn{5} Transverse waves are set in a thin rod of length $\ell$ clamped at one end O and free at other end. The wave velocity is $v$. The rod vibrates in the 4$^{\text{th}}$ overtone mode with amplitude $y_0$. The equation of motion of the point at $x=\dfrac{2\ell}3$ is
\qfig{p1_phy_s1_q05_1}
\begin{qoptions}
\opt{A}{$y=y_0\sin\dfrac{5\pi}3\cos\dfrac{5\pi vt}{2\ell}$}
\opt{B}{$y=y_0\sin\dfrac{7\pi}3\cos\dfrac{7\pi vt}{2\ell}$}
\opt{C}{$y=y_0\sin10\dfrac\pi3\cos\dfrac{9\pi vt}{2\ell}$}
\opt{D}{$y=0$}
\end{qoptions}

\qn{6} Four wire A, B, C and D each of length $\ell=10$cm and each of area of cross section is 0.1 m$^2$ are connected in the given circuit. Then, the position of null point is :- (Given that resistivity : $\rho_A=1\Omega$ - m, $\rho_B=3\Omega$ - m, $\rho_C=6\Omega$ - m, $\rho_D=1\Omega$ - m)
\qfig{p1_phy_s1_q06_1}
\begin{qoptions}
\opt{A}{mid point of wire B or wire C when both the switches S$_1$ and S$_2$ open.}
\opt{B}{mid point of wire B when both the switches S$_1$ and S$_2$ are closed.}
\opt{C}{mid point of wire D when both the switches S$_1$ and S$_2$ are open.}
\opt{D}{None of these}
\end{qoptions}

\qn{7} Uniform, smooth, thin, fairly long rod of mass $M$ and length $L$ is in weightlessness in the distant space as shown. A small bead is placed on the rod, the mass $m$ of which is much less than the mass $M$ of the rod ($m\ll M$). Find the period $T$ of small oscillations of the bead near the center rod. Gravitational constant is $G$.
\qfig{p1_phy_s1_q07_1}
\begin{qoptions}
\opt{A}{$\pi\sqrt{\dfrac{L^3}{8GM}}$}
\opt{B}{$\pi\sqrt{\dfrac{L^3}{GM}}$}
\opt{C}{$2\pi\sqrt{\dfrac{L^3}{2GM}}$}
\opt{D}{$2\pi\sqrt{\dfrac{L^3}{8GM}}$}
\end{qoptions}

\qn{8} Consider three identical uniform disc each of mass $M$ and radius `$R$'. System is kept in $x$-$y$ plane. Match the moment of inertia about axis in List I to answer in List II :-
\qfig{p1_phy_s1_q08_1}
\par\nobreak
\begin{center}\small
\begin{tabular}{|c|p{2.4cm}|c|p{2.6cm}|}\hline
 & \multicolumn{1}{c|}{\textbf{List$-$I}} & & \multicolumn{1}{c|}{\textbf{List$-$II}}\\\hline
(I) & About $x$ axis & (P) & $\frac{18+4\sqrt3}4MR^2$\\\hline
(II) & About $y$ axis & (Q) & $\dfrac{11MR^2}4$\\\hline
(III) & About an axis passing through centre of mass of system and perpendicular to its plane & (R) & $\dfrac{11MR^2}2$\\\hline
(IV) & About axis passing through point A and perpendicular to plane & (S) & $\frac{27+8\sqrt3}4MR^2$\\\hline
 & & (T) & None of these\\\hline
\end{tabular}
\end{center}
\begin{qoptions}
\opt{A}{(I) - (S) ; (II) - (Q) ; (III) - (R) ; (IV) - (T)}
\opt{B}{(I) - (T) ; (II) - (Q) ; (III) - (R) ; (IV) - (S)}
\opt{C}{(I) - (S) ; (II) - (R) ; (III) - (Q) ; (IV) - (T)}
\opt{D}{(I) - (P) ; (II) - (Q) ; (III) - (S) ; (IV) - (R)}
\end{qoptions}

\qn{9} A wire of length $L$, mass $m$ and carrying a current $I$ is suspended from point O as shown. An another infinitely long wire carrying the same current $I$ is at a distance $L$ below the lower end of the wire. The angular acceleration of the wire just after it is released from the position is :-\\
(given $I=2$ A, $L=1$m and $m=0.1$ kg, $\log_e2=0.7$)
\qfig{p1_phy_s1_q09_1}
\begin{qoptions}
\opt{A}{$48\times10^{-7}$ rad/s$^2$}
\opt{B}{$96\times10^{-7}$ rad/s$^2$}
\opt{C}{$24\times10^{-6}$ rad/s$^2$}
\opt{D}{$48\times10^6$ rad/s$^2$}
\end{qoptions}

\qn{10} 2 identical cylinder-shaped barrels, having the same height $h$ are placed on the horizontal ground, one on the top of the other. Both were nearly fully filled with wine. Two holes are drilled as shown. Find $x$, in order that the impact points of the wine rays fall as far apart as possible.
\qfig{p1_phy_s1_q10_1}
\begin{qoptions}
\opt{A}{$\dfrac12$}
\opt{B}{$\dfrac13$}
\opt{C}{$\dfrac14$}
\opt{D}{$\dfrac23$}
\end{qoptions}

\qn{11} Three identical large parallel plates are fixed at separation $d$ from each other as shown. The area of each plate is $A$. Plate 1 is given charge $+Q$, while plates 2 and 3 are neutral and are connected to each other through coil of inductance $L$ and switch S. If resistance of all connecting wires is neglected then the maximum current that will flow through coil after closing switch is\\
(take $C=\dfrac{\varepsilon_0A}d$ and neglect fringe effect)
\qfig{p1_phy_s1_q11_1}
\begin{qoptions}
\opt{A}{$\dfrac Q{\sqrt{LC}}$}
\opt{B}{$\dfrac Q{\sqrt{2LC}}$}
\opt{C}{$\dfrac{3Q}{2\sqrt{LC}}$}
\opt{D}{$\dfrac Q{2\sqrt{LC}}$}
\end{qoptions}

\qn{12} Which of the following statements are correct regarding L.E.D.?\\
(I) Low operational voltage and less power.\\
(II) Fast action and no warm-up time required.\\
(III) The bandwidth of emitted light is 100 \AA\ to 500 \AA\ or in other words it is nearly (but not exactly) monochromatic.\\
(IV) Short life and ruggedness.\\
(V) Fast on-off switching capacity.
\begin{qoptions}
\opt{A}{I, II, III, IV are correct.}
\opt{B}{I, III, IV, V are correct}
\opt{C}{I, II, III, V are correct}
\opt{D}{All are correct.}
\end{qoptions}

\qn{13} A stone is thrown vertically upward. On its way up it passes point A with speed of $v$, and point B, 3m higher than A, with speed of $v/2$. The maximum height reached by stone above point B is :
\begin{qoptions}
\opt{A}{1m}
\opt{B}{2m}
\opt{C}{3m}
\opt{D}{5m}
\end{qoptions}

\qn{14} A particle of mass $m$ is driven by a machine that delivers a constant power $k$ watts. If the particle starts from rest the force on the particle at time $t$ is :
\begin{qoptions}
\opt{A}{$\sqrt{2mk}\ t^{-1/2}$}
\opt{B}{$\dfrac12\sqrt{mk}\ t^{-1/2}$}
\opt{C}{$\sqrt{\dfrac{mk}2}\ t^{-1/2}$}
\opt{D}{$\sqrt{mk}\ t^{-1/2}$}
\end{qoptions}

\qn{15} A solid sphere of mass $m$ and radius $R$ is gently placed on a conveyer belt moving with constant velocity $v_0$. If coefficient of friction between belt and sphere is $2/7$, the distance travelled by the centre of the sphere before it starts pure rolling is :
\qfig{p1_phy_s1_q15_1}
\begin{qoptions}
\opt{A}{$\dfrac{v_0^2}{7g}$}
\opt{B}{$\dfrac{2v_0^2}{49g}$}
\opt{C}{$\dfrac{2v_0^2}{5g}$}
\opt{D}{$\dfrac{2v_0^2}{7g}$}
\end{qoptions}

\qn{16} An object of mass 0.2 kg executes simple harmonic oscillations along X-axis with a frequency of $(25/\pi)$ Hz. At the position $x=0.04$ m, the object has kinetic energy of 0.5 J and potential energy of 0.4 J. Find the amplitude of oscillations (in cm). [assume zero potential energy at equilibrium position]
\begin{qoptions}
\opt{A}{2 cm}
\opt{B}{4 cm}
\opt{C}{5 cm}
\opt{D}{6 cm}
\end{qoptions}

\qn{17} The de-Broglie wavelength associated with electron in $n=4$ level is :
\begin{qoptions}
\opt{A}{two times the de-Broglie wavelength of the electron in the ground state}
\opt{B}{four times the de-Broglie wavelength of the electron in the ground state}
\opt{C}{half of the de-Broglie wavelength of the electron in the ground state}
\opt{D}{$1/4^{\text{th}}$ of the de-Broglie wavelength of the electron in the ground state.}
\end{qoptions}

\qn{18} A mixture of light, consisting of wavelength 590 nm and an unknown wavelength, illuminates Young's double slit and gives rise to two overlapping interference patterns on the screen. The central maximum of both lights coincide. Further, it is observed that the third bright fringe of known light coincides with the 4th bright fringe of the unknown light. From this data, the wavelength of the unknown light is :
\begin{qoptions}
\opt{A}{393.4 nm}
\opt{B}{885.0 nm}
\opt{C}{442.5 nm}
\opt{D}{776.8 nm}
\end{qoptions}

\qn{19} An organ pipe of cross sectional area 100 cm$^2$ resonates with a tuning fork of frequency 1000Hz in fundamental tone. The minimum volume of water to be drained out so that the pipe again resonates with the same tuning fork is\\
(take velocity of wave = 320 m/s)
\begin{qoptions}
\opt{A}{800 cm$^3$}
\opt{B}{1200cm$^3$}
\opt{C}{1600cm$^3$}
\opt{D}{2000 cm$^3$}
\end{qoptions}

\qn{20} The electric field potential in a certain region of space depends only on the $x$-coordinate as $V=-2x^3+3$. Find the distribution of the space charge $\rho$ ($x$).
\begin{qoptions}
\opt{A}{$6\in_0x$}
\opt{B}{$12\in_0x$}
\opt{C}{$6\in_0$}
\opt{D}{$-12\in_0$}
\end{qoptions}

\sect{Section-II}

\qn{1} As shown, the left end of an upright U-shaped glass tube of uniform cross section is sealed and the right end is open. If the atmospheric pressure is $P_0=76$ cm Hg and at $T_0=15^\circ$C, the length of the gas column sealed by the mercury is $\ell_0=8$ cm and the height difference between the left and the right mercury surfaces is $h=4$ cm. When the temperature is raised to $T_1$, the length of the gas column becomes $\ell_1=9$cm, then $T_1$ is $n\times12^\circ$C. Then find value of $n$.
\qfig{p1_phy_s2_q01_1}

\qn{2} An infinitely long wire, located on the $Z$-axis, carries a current $I$ along the $+Z$-direction and produces the magnetic field $\vec B$. The magnitude of the line integral $\displaystyle\int\vec B\cdot\vec{dr}$ along a straight line from the point $\left(-\sqrt3a,a,0\right)$ to $(a,a,0)$ is given by $\dfrac{x\mu_0I}{24}$ then find $x$. [$\mu_0$ is the magnetic permeability of free space]

\qn{3} A length $L=15$m of flexible tape is tightly wound. It is then allowed to unwinding as it rolls down a fixed incline that makes an angle $\theta=30^\circ$ with the horizontal, the upper end of the tape being fixed. Find the time taken (in second) by the tape to unwind completely. Neglect radius at any time w.r.t. the length of the rope.
\qfig{p1_phy_s2_q03_1}

\qn{4} The number of photons of yellow light ($\lambda=575$ nm) that must be absorbed by 1.0 cm$^2$ of human skin to raise the skin temperature by 1.0$^\circ$C is $1.2\times10^n$. Find $n$. (Assume that the photon energy is absorbed by 1.0 cm$^3$ of tissue. As rough estimates of the tissue's density and specific heat, use the values for water)

\qn{5} A projectile is fired at a speed of 100 m/s at an angle of 37$^\circ$ above the horizontal. At the highest point, the projectile breaks into two parts of mass ratio 1 : 3. The smaller part comes to the rest. Find the distance (in m) from the launching point to the point where the heavier part lands [$g=10\ \mathrm{m/s^2}$]

\partsec{2}{CHEMISTRY}
\sect{Section-I}

\qn{1} During the qualitative analysis of bromide ion by chlorine water, orange brown coloration is obtained in organic layer due to disolution of bromine.\\
Which organic compound is used in above test to prepare organic layer.\\
(P) CCl$_4$\\
(Q) CS$_2$\\
(R) CH$_2$Cl$_2$\\
(S) C$_2$H$_5$OH
\begin{qoptions}
\opt{A}{(P), (Q) and (S)}
\opt{B}{(P) , (Q) only}
\opt{C}{(P), (Q), (R), (S)}
\opt{D}{(P), (Q) and (R)}
\end{qoptions}

\qn{2} For an ideal gas, three adiabatic processes are carried out upto same final pressure from same initial state. If adiabatic reversible process ends up at point `B' and adiabatic single step irreversible process ends up at point `C', then adiabatic free expansion upto same final pressure will end up at -
\qfig{p1_chem_s1_q02_1}
\begin{qoptions}
\opt{A}{A}
\opt{B}{P}
\opt{C}{Q}
\opt{D}{R}
\end{qoptions}

\qn{3}\par\nobreak
\qfig{p1_chem_s1_q03_1}
Calculate total moles of CO$_2$ evolved per mol substrate
\begin{qoptions}
\opt{A}{5}
\opt{B}{4}
\opt{C}{6}
\opt{D}{7}
\end{qoptions}

\qn{4} When hydrated ferric chloride is dissolved in aqueous solution of oxalic acid containing potassium hydroxide
\begin{qoptions}
\opt{A}{Brown ppt of Fe(OH)$_3$ is obtained}
\opt{B}{Green crystals of K$_3$[Fe(C$_2$O$_4$)$_3$] are obtained}
\opt{C}{Yellow crystals of K[FeCl$_4$] are obtained}
\opt{D}{Yellow solution of [Fe(H$_2$O)$_6$]Cl$_3$ \& K$_2$C$_2$O$_4$ are obtained}
\end{qoptions}

\qn{5} For a given ideal gas, the mean free path\\
(at a particular pressure) is -
\begin{qoptions}
\opt{A}{Independent of temperature}
\opt{B}{Decreases with rise in temperature}
\opt{C}{Increases with rise in temperature}
\opt{D}{Directly proportional to T$^2$}
\end{qoptions}

\qn{6}\par\nobreak
\qfig{p1_chem_s1_q06_1}
\begin{qoptions}
\opt{A}{C cannot exhibit GI}
\opt{B}{\qfigi{p1_chem_s1_q06_2}}
\opt{C}{\qfigi{p1_chem_s1_q06_4}}
\opt{D}{\qfigi{p1_chem_s1_q06_3}}
\end{qoptions}

\qn{7} For a gaseous reaction :\\
CH$_3$CHO(g) $\xrightarrow{\ \Delta\ }$ CO(g) + CH$_4$(g)\\
initial pressure is 80 mm of Hg and total pressure at the end of 20 minutes is 120 mm of Hg. The rate constant of the reaction assuming first order kinetics is - (ln2 = 0.693)
\begin{qoptions}
\opt{A}{$3.465\times10^{-2}$ min$^{-1}$}
\opt{B}{34.65 min$^{-1}$}
\opt{C}{3.465 min$^{-1}$}
\opt{D}{0.3465 min$^{-1}$}
\end{qoptions}

\qn{8} \textbf{Assertion (A) :} Actinoid contraction is greater from element to element than lanthanide contraction.\\
\textbf{Reason (R) :} The 5f electrons themselves provide poor shielding from element to element in series.
\begin{qoptions}
\opt{A}{A is false but R is true.}
\opt{B}{A is true but R is false.}
\opt{C}{Both A and R are true and R is the correct explanation of A.}
\opt{D}{Both A and R are true and R is NOT the correct explanation of A.}
\end{qoptions}

\qn{9} pH of the aq. solution of 0.05 M Ca(CN)$_2$ at 25$^\circ$C is :\\
$\left\{\mathrm{pK_b(CN^-)=8.26}\right\}$
\begin{qoptions}
\opt{A}{10.63}
\opt{B}{9.37}
\opt{C}{8.63}
\opt{D}{8.37}
\end{qoptions}

\qn{10} Which of the following compounds has least equilibrium constant for acid catalyzed hydration :
\begin{qoptions}
\opt{A}{\qfigi{p1_chem_s1_q10_2}}
\opt{B}{\qfigi{p1_chem_s1_q10_1}}
\opt{C}{\qfigi{p1_chem_s1_q10_3}}
\opt{D}{\qfigi{p1_chem_s1_q10_4}}
\end{qoptions}

\qn{11} In the laboratory a aqueous manganese (II) ion salts is oxidised by \underline{\hspace{2em}} to \underline{\hspace{2em}}.\\
Select best option to fill the blank space respectively.
\begin{qoptions}
\opt{A}{Sulphate ; Manganate}
\opt{B}{Peroxydisulphate ; Manganate}
\opt{C}{Thiosulphate ; Permanganate}
\opt{D}{Peroxydisulphate ; Permanganate}
\end{qoptions}

\qn{12} Compound PdCl$_4$.6H$_2$O is a hydrated complex. If 1 molal aqueous solution of this complex has freezing point of 269.28 K (assuming 100\% ionization of this complex), then formula of the complex is :\\
(K$_f$ for water = 1.86 K kg mol$^{-1}$)
\begin{qoptions}
\opt{A}{$\left[\mathrm{Pd(H_2O)_6}\right]\mathrm{Cl_4}$}
\opt{B}{$\left[\mathrm{Pd(H_2O)_4Cl_2}\right]\mathrm{Cl_2.2H_2O}$}
\opt{C}{$\left[\mathrm{Pd(H_2O)_3Cl_3}\right]\mathrm{Cl.3H_2O}$}
\opt{D}{$\left[\mathrm{Pd(H_2O)_2Cl_4}\right].4\mathrm{H_2O}$}
\end{qoptions}

\qn{13}\par\nobreak
\qfig{p1_chem_s1_q13_1}
Select the correct statement about final nitrogen containing compound [D].
\begin{qoptions}
\opt{A}{It is less basic than aniline.}
\opt{B}{Formal charge of N is +1.}
\opt{C}{It gives bad smelling compound on treatment with CHCl$_3$/KOH.}
\opt{D}{On treatment with NaNO$_2$ + HCl followed by $\beta-$naphthol in alkaline medium, it forms red dye.}
\end{qoptions}

\qn{14} Consider Be, B, C, N, O, F, Cl, S atoms and identify atom which has maximum E.A. and atom which has lowest ionization energy.\\
What is the shape of binary molecule formed by them.
\begin{qoptions}
\opt{A}{Linear}
\opt{B}{Tetrahedral}
\opt{C}{T-pyramidal}
\opt{D}{Triangular planar}
\end{qoptions}

\qn{15}\par\nobreak
\qfig{p1_chem_s1_q15_1}
Choose the correct statement regarding `R' and `P'.
\begin{qoptions}
\opt{A}{When R = dil.HNO$_3$, number of Stereoisomers of P = 10}
\opt{B}{When R = Br$_2$ / H$_2$O, number of enantiomeric pair of P = 6}
\opt{C}{When R = Red P + HI, number of structural isomers of P = 6}
\opt{D}{When R = HIO$_4$ (5 moles), number of moles of acid product obtained per mole substrate is 6.}
\end{qoptions}

\qn{16} Select the p-block non metallic atoms present in chalcopyrite, cryolite and pyrolusite respectively.
\begin{qoptions}
\opt{A}{Sulphur, Fluorine, Oxygen}
\opt{B}{Iron, Aluminium, Manganese}
\opt{C}{Sulphur, Fluorine, Sulphur}
\opt{D}{Sulphur, Chlorine, Fluorine}
\end{qoptions}

\qn{17} In the following reaction sequence\\
CH$_3-$CH$_2-$Cl $\xrightarrow{\ \text{KCN}\ }\ \xrightarrow{\ \text{C}_6\text{H}_6\text{MgBr}\ }\ \xrightarrow{\ \text{H}_3\text{O}^{+}\ }$ P\\
The IUPAC name of `P' is.
\begin{qoptions}
\opt{A}{1-Phenylpropanone}
\opt{B}{1-Phenylpropanoic acid}
\opt{C}{1-Phenylethanal}
\opt{D}{1-Phenylethanoic acid}
\end{qoptions}

\qn{18} Select the set in which atleast two complexes are heteroleptic and contain chelate rings.
\begin{qoptions}
\opt{A}{Set-1\\ Pentaammineaquacobalt(III) chloride\\ Tris(ethylenediamine)cobalt(III) nitrate\\ Sodium trioxalatoaluminate(III)}
\opt{B}{Set-2\\ Diamminebis(ethylenediamine)cobalt(III) chloride\\ Bis(ethylenediamine)difluoridocobalt(III) nitrate\\ Sodium trioxalatoaluminate(III)}
\opt{C}{Set-3\\ Pentaammineaquacobalt(III) chloride\\ Bis(ethylenediamine)difluoridocobalt(III) nitrate\\ Sodium trioxalatoaluminate(III)}
\opt{D}{Set-4\\ Diamminebis(ethylenediamine)cobalt(III) chloride\\ Pentaammineaquacobalt(III) chloride\\ Tris(ethylenediamine)cobalt(III) nitrate}
\end{qoptions}

\qn{19} \textbf{Statement -1 :} The SRP of three metallic ions A$^{+}$, B$^{2+}$, C$^{3+}$ are $-0.3$, $-0.5$, 0.8 volt respectively, so oxidising power of ions is C$^{3+}$ $>$ A$^{+}$ $>$ B$^{2+}$.\\
\textbf{Statement -2 :} Higher the SRP, higher the oxidising power.
\begin{qoptions}
\opt{A}{Statement-1 is true, statement-2 is true and statement-2 is correct explanation for statement-1.}
\opt{B}{Statement-1 is true, statement-2 is true and statement-2 is NOT the correct explanation for statement-1.}
\opt{C}{Statement-1 is true, statement-2 is false.}
\opt{D}{Statement-1 is false, statement-2 is true.}
\end{qoptions}

\qn{20} Ph$-$N$_2^{+}$Cl$^{-}$ is made to react with various reagents as follows
\par\nobreak
\begin{center}\small
\begin{tabular}{|c|p{2.4cm}|c|p{2.6cm}|}\hline
\multicolumn{2}{|c|}{\shortstack{\textbf{List-I}\\\textbf{(Reagent)}}} & \multicolumn{2}{c|}{\shortstack{\textbf{List-II}\\\textbf{(Product)}}}\\\hline
(P) & \qfigi{p1_chem_s1_q20_1} & (1) & \qfigi{p1_chem_s1_q20_2}\\\hline
(Q) & HBF$_4$, $\Delta$ & (2) & \qfigi{p1_chem_s1_q20_3}\\\hline
(R) & Cu, HCl & (3) & \qfigi{p1_chem_s1_q20_4}\\\hline
(S) & CuCN/KCN & (4) & \qfigi{p1_chem_s1_q20_5}\\\hline
\end{tabular}
\end{center}
Match List I with List II and choose the correct answer from the options given below:
\begin{qoptions}
\opt{A}{P$\to$4, Q$\to$3, R$\to$2, S$\to$1}
\opt{B}{P$\to$1, Q$\to$3, R$\to$4, S$\to$2}
\opt{C}{P$\to$3, Q$\to$1, R$\to$2, S$\to$4}
\opt{D}{P$\to$3, Q$\to$1, R$\to$4, S$\to$2}
\end{qoptions}

\sect{Section-II}

\qn{1} Find sum of the bond order of species which must have atleast one $\pi$ bond among :-\\
H$_2$ , N$_2^{+}$ , O$_2^{+}$ , C$_2^{2-}$ , F$_2$ , C$_2$

\qn{2} Wavelength of electron waves in two Bohr orbits of a H-like species is in ratio 3 : 5. The ratio of kinetic energy of electron in these two Bohr orbits is 25 : $x$, hence $x$ is :

\qn{3} Calculate enthalpy of precipitation of CaCO$_3$(s) (in kcal mol$^{-1}$).\\
Ca$^{2+}$(aq.) + CO$_3^{2-}$ (aq.) $\to$ CaCO$_3$ (s)\\
Given : $\Delta_f$H$^0$(Ca$^{2+}$, aq.) $=-130$ kcal mol$^{-1}$\\
\hspace*{3.2em}$\Delta_f$H$^0$(CO$_3^{2-}$, aq.) $=-160$ kcal mol$^{-1}$\\
\hspace*{3.2em}$\Delta_f$H$^0$(CaCO$_3$, s) $=-285$ kcal mol$^{-1}$

\qn{4} Total number of enols for the product formed in the given reaction will be (Including stereoisomers)
\qfig{p1_chem_s2_q04_1}

\qn{5} How many compounds show geometrical isomerism after reaction with NH$_2-$OH?
\qfig{p1_chem_s2_q05_1}

\partsec{3}{MATHEMATICS}
\sect{Section-I}

\qn{1} Let A be the set of all real solutions of equation $x\left(x^2+3|x|+5|x-1|+6|x-2|\right)=0$ and B be the set of all real solutions of equation $x^2-|x|-12=0$ then number of subsets of the set $A\times B$ is
\begin{qoptions}
\opt{A}{2}
\opt{B}{4}
\opt{C}{8}
\opt{D}{16}
\end{qoptions}

\qn{2} Let $\alpha,\beta,\gamma$ be three roots of equation $x^3+bx+c=0$ such that $\beta\gamma=1=-\alpha$. If $\alpha^n+\beta^n+\gamma^n=-3$ then value of $n$ can be
\begin{qoptions}
\opt{A}{9}
\opt{B}{8}
\opt{C}{5}
\opt{D}{6}
\end{qoptions}

\qn{3} \textbf{Assertion A :} The number of ways in which 3 married couples with their 4 children can sit in a row such that no husband and wife are together, are $10!-3.9!+3.8!-7!$.\\
\textbf{Reason R :} Number of ways of occurrence of at least one event out of three events A, B \& C $=n(A\cup B\cup C)=n(A)+n(B)+n(C)-n(A\cap B)-n(B\cap C)-n(C\cap A)+n(A\cap B\cap C)$.
\begin{qoptions}
\opt{A}{Both \textbf{A} and \textbf{R} are correct but \textbf{R} is \textbf{NOT} the correct explanation of \textbf{A}}
\opt{B}{Both \textbf{A} and \textbf{R} are correct and \textbf{R} is the correct explanation of \textbf{A}}
\opt{C}{\textbf{A} is correct but \textbf{R} is not correct}
\opt{D}{\textbf{A} is not correct but \textbf{R} is correct}
\end{qoptions}

\qn{4} If for the complex number $z$ satisfying $|z-2-2i|\leq1$, the maximum and minimum value of $|3iz+6|$ is $M$ and $m$ then value of $(M+m)$ is
\begin{qoptions}
\opt{A}{9}
\opt{B}{3}
\opt{C}{12}
\opt{D}{15}
\end{qoptions}

\qn{5} If $2\tan^2\theta-5\sec\theta=1$ has exactly 7 solutions in the interval $\left[0,\dfrac{n\pi}2\right]$, then for the least value of $n\in N$ value of $\displaystyle\sum_{k=1}^n\dfrac k{2^k}$ is equal to
\begin{qoptions}
\opt{A}{$\dfrac1{2^{13}}\left(2^{14}-15\right)$}
\opt{B}{$1-\dfrac{15}{2^{13}}$}
\opt{C}{$\dfrac1{2^{15}}\left(2^{14}-14\right)$}
\opt{D}{$\dfrac1{2^{14}}\left(2^{15}-15\right)$}
\end{qoptions}

\qn{6} Let the first term of the series be $a_1=6$ and its $r^{\text{th}}$ term be $a_r=3.a_{r-1}+6^r$, $r=2,3,\dots n$. For some $n\in N$, if the sum of first $n$ terms of the series is $\dfrac15\left(n^2-18n+84\right)\left(4.6^n-5.3^n+1\right)$ then $n$ is equal to
\begin{qoptions}
\opt{A}{9}
\opt{B}{6}
\opt{C}{7}
\opt{D}{8}
\end{qoptions}

\qn{7} Consider the system of linear equations\\
$2x+\left(p^2-2\right)y+6z=8$\\
$x+2y+(q-1)z=5$\\
$x+y+3z=4$ where $p,q\in R$\\
which of the following statements is NOT TRUE
\begin{qoptions}
\opt{A}{System will have unique solution if $p\in R-\{-2,2\}$, $q\in R-\{4\}$}
\opt{B}{system is inconsistent if $p=2$, $q=4$}
\opt{C}{system will have infinite solution if $p=-2$, $q\in R$}
\opt{D}{system is consistent if $p\in R$, $q\in R-\{4\}$}
\end{qoptions}

\qn{8} The value of $\displaystyle\int_{-\alpha}^{-1}\left(\sec^{-1}y-\tan^{-1}\sqrt{y^2-1}\right)dy+\int_1^\alpha\left(\sec^{-1}y-\tan^{-1}\sqrt{y^2-1}\right)dy$ is :\\
(where $\alpha>1$ and $\displaystyle\int_1^\alpha\sec^{-1}(y)dy=\beta$)
\begin{qoptions}
\opt{A}{$\pi(\alpha-1)-\beta$}
\opt{B}{$\pi\alpha-2\beta$}
\opt{C}{$\pi\alpha-2(\beta-1)$}
\opt{D}{$\pi(\alpha-1)-2\beta$}
\end{qoptions}

\qn{9} Match List-I with List-II and select the correct answer using the code given below the list.
\par\nobreak
\begin{center}\small
\begin{tabular}{|c|p{3.9cm}|c|c|}\hline
 & \multicolumn{1}{c|}{\textbf{List$-$I}} & & \textbf{List$-$II}\\\hline
(A) & Let $\left(\vec a\times\vec b\right)\times\vec c=-5\vec a+4\vec b$ and $\vec a.\vec b=3$ and $\vec a\times\left(\vec b\times\vec c\right)=\lambda\vec b+\mu\vec c$ ($\lambda,\mu$ are scalars), then $|\lambda|+|\mu|$ is (Where $\vec a,\vec b,\vec c$ are pairwise non-collinear) & (I) & 2\\\hline
(B) & If the line $y-\sqrt3x+3=0$ cuts the parabola $y^2=x+2$ at A and B, given that $|(PA)(PB)|=\dfrac{\lambda_1\left(2+\sqrt3\right)}{\lambda_2}$, then $\dfrac{\lambda_1^2+\lambda_2^2}5$ is (where P is $\left(\sqrt3,0\right)$ and $\lambda_1,\lambda_2$ are relatively prime) & (II) & 7\\\hline
(C) & The sum (in radians) of all values of $x$ with $0\leq x\leq2\pi$ which satisfy $\sqrt2(\cos2x-\sin x-1)=1+2\sin x$, is $k\pi$, then $k$ is & (III) & 6\\\hline
(D) & Value of $\displaystyle\int_{-\pi/4}^{\pi/4}\dfrac{x^9-3x^7+x^5+1}{\cos^2x}dx$ is equal to & (IV) & 5\\\hline
\end{tabular}
\end{center}
\begin{qoptions}
\opt{A}{(A) $-$ I, (B) $-$ II, (C) $-$ III, (D) $-$ IV}
\opt{B}{(A) $-$ II, (B) $-$ IV, (C) $-$ III, (D) $-$ I}
\opt{C}{(A) $-$ IV, (B) $-$ I, (C) $-$ III, (D) $-$ II}
\opt{D}{(A) $-$ III, (B) $-$ IV, (C) $-$ II, (D) $-$ I}
\end{qoptions}

\qn{10} Let $y=y(x)$ be the solution of the differential equation $(\text{cosec}\,x)dy+(2(1-x)\cot x+x(x-2))dx=0$ such that $y\left(\dfrac\pi2\right)=3$. The value of $y(2)$ is equal to
\begin{qoptions}
\opt{A}{2}
\opt{B}{$2(1-\cos2)$}
\opt{C}{$2(1+\cos2)$}
\opt{D}{3}
\end{qoptions}

\qn{11} Consider the given data frequency distribution :
\begin{center}\small
\begin{tabular}{|l|c|c|c|c|c|}\hline
Class interval & 0-6 & 6-12 & 12-18 & 18-24 & 24-30\\\hline
Frequency & 4 & 5 & 3 & 6 & 2\\\hline
\end{tabular}
\end{center}
The mean deviation from median of given data is
\begin{qoptions}
\opt{A}{7}
\opt{B}{7.5}
\opt{C}{6}
\opt{D}{6.5}
\end{qoptions}

\qn{12} Given a real valued function $f(x)$ defined as\\
$f(x)=\begin{cases}\dfrac{\tan^2\{x\}}{\left(\{x\}\right)^2} & ;\ x>0\\ 1 & ;\ x=0\\ \sqrt{\{x\}\cot\{x\}} & ;\ x<0\end{cases}$\\
where $\{x\}$ represent fractional part of $x$. Then which of the following is INCORRECT?
\begin{qoptions}
\opt{A}{$\displaystyle\lim_{x\to0^+}f(x)=1$}
\opt{B}{$\cot^{-1}\left(\displaystyle\lim_{x\to0^-}f(x)\right)^2=1$}
\opt{C}{$\tan^{-1}\left(\displaystyle\lim_{x\to0^+}f(x)\right)=\dfrac\pi4$}
\opt{D}{$\displaystyle\lim_{x\to0}f(x)$ exists.}
\end{qoptions}

\qn{13} A square is inscribed in the circle $x^2+y^2-10x-6y+30=0$. One side of this square is parallel to $y=x+3$. If $(x_i,y_i)$ (where $i=1,2,3,4$) are the vertices of square then $\displaystyle\sum_{i=1}^4\left(x_i^2+y_i^2\right)$ is equals to
\begin{qoptions}
\opt{A}{148}
\opt{B}{156}
\opt{C}{160}
\opt{D}{152}
\end{qoptions}

\qn{14} Let P be the point $(10,-2,-1)$ and Q be the foot of the perpendicular drawn from the point R$(1,7,6)$ on the line passing through the points $(2,-5,11)$ and $(-6,7,-5)$. Then the length of the line segment PQ is equal to
\begin{qoptions}
\opt{A}{14}
\opt{B}{13}
\opt{C}{12}
\opt{D}{17}
\end{qoptions}

\qn{15} \textbf{Statement-1 :} If $a+b+c+d=9$ and $a^2+b^2+c^2+d^2=27$, where $a,b,c,d$ are non-negative real numbers, then $d\in\left[0,\dfrac92\right]$.\\
\textbf{Statement-2 :} $\left(\dfrac{a+b+c}3\right)^2\leqslant\dfrac{a^2+b^2+c^2}3$.
\begin{qoptions}
\opt{A}{Statement-1 is True, Statement-2 is True}
\opt{B}{Statement-1 is False, Statement-2 is False}
\opt{C}{Statement-1 is True, Statement-2 is False}
\opt{D}{Statement-1 is False, Statement-2 is True}
\end{qoptions}

\qn{16} Let unit vector $\vec c$ is inclined at an angle $\theta$ to a unit vector $\vec a$ and $\vec b$ which are perpendicular to each other. If $\vec c=\lambda\left(\vec a+\vec b\right)+\mu\left(\vec a\times\vec b\right)$ where $\lambda$, $\mu$ be real, then $\theta$ belongs to
\begin{qoptions}
\opt{A}{$\left(-\dfrac\pi4,0\right)$}
\opt{B}{$\left[0,\dfrac\pi4\right]$}
\opt{C}{$\left[\dfrac{3\pi}4,\pi\right)$}
\opt{D}{$\left[\dfrac\pi4,\dfrac{3\pi}4\right]$}
\end{qoptions}

\qn{17} Length of focal chord of parabola $y^2=8x$ parallel to normal drawn to it at point with abscissa 2 is -
\begin{qoptions}
\opt{A}{4}
\opt{B}{8}
\opt{C}{12}
\opt{D}{16}
\end{qoptions}

\qn{18} Let P be a relation defined on the set of interval $\left(0,\dfrac\pi2\right]$ such that $P=\{(a,b):\text{cosec}^2a-\cot^2b=1\}$. Then P is
\begin{qoptions}
\opt{A}{Reflexive and symmetric but not transitive}
\opt{B}{Reflexive and transitive but not symmetric}
\opt{C}{Symmetric and transitive but not reflexive}
\opt{D}{Equivalence relation}
\end{qoptions}

\qn{19} The area enclosed by $y=\ell nx$, $y=\ell n|x|$, $y=|\ell nx|$ and $y=|\ell n|x||$ is
\begin{qoptions}
\opt{A}{5}
\opt{B}{2}
\opt{C}{4}
\opt{D}{1}
\end{qoptions}

\qn{20} Let $J=\displaystyle\int_{-5}^{-4}\left(3-x^2\right)\tan\left(3-x^2\right)dx$\\
$K=\displaystyle\int_{-2}^{-1}\left(6-6x+x^2\right)\tan\left(6x-x^2-6\right)dx$.\\
Then $(J+K)$
\begin{qoptions}
\opt{A}{0}
\opt{B}{1}
\opt{C}{2}
\opt{D}{3}
\end{qoptions}

\sect{Section-II}

\qn{1} The solution of the equation $\dfrac8{\{x\}}=\dfrac9x+\dfrac{10}{[x]}$ is of the form $\dfrac{k+1}k$, $k\in N$ then $k$ is equal to\\
(where $[x]$ denotes greatest integer less than $x$ \& $\{x\}$ denotes fractional part of $x$)

\qn{2} If $\displaystyle\sum_{n=0}^\infty2\cot^{-1}\left(\dfrac{n^2+n+4}2\right)=k\pi$, then value of $k$ is

\qn{3} If $a$ and $b$ are chosen randomly from the set $A=\{1,2,3,4,5,6\}$ with replacement. If probability that $\displaystyle\lim_{x\to0}\left(\dfrac{a^x+b^x}2\right)^{\frac2x}=6$ is $\dfrac pq$ [H.C.F(p,q) = 1] then $(q-p)$ is

\qn{4} If $A=\begin{bmatrix}1&2\\3&4\end{bmatrix}$ and $A^{2014}=\lambda A^{2013}+\mu A^{2012}$ then $(\lambda+\mu)$ is

\qn{5} The coefficient of $x^{\frac{n^2+n-14}2}$ in\\
$(x-1)\left(x^2-2\right)\left(x^3-3\right)\dots\left(x^n-n\right)$, $n\geqslant30$ is

\end{multicols}
\end{document}
